% Generated draft: requires mathematical and source review.
\documentclass[11pt]{article}
\usepackage[margin=25mm]{geometry}
\usepackage{fontspec}
\setmainfont{Latin Modern Roman}
\usepackage{amsmath,amssymb,mathtools,unicode-math}
\setmathfont{Latin Modern Math}
\usepackage{enumitem,needspace,xurl,hyperref}
\setlength{\parindent}{0pt}
\setlength{\parskip}{5pt}
\setlength{\emergencystretch}{3em}
\newcommand{\N}{\mathbb{N}}
\newcommand{\Z}{\mathbb{Z}}
\newcommand{\Q}{\mathbb{Q}}
\newcommand{\R}{\mathbb{R}}
\newcommand{\C}{\mathbb{C}}
\newcommand{\F}{\mathbb{F}}
\newcommand{\A}{\mathbb{A}}
\newcommand{\PP}{\mathbb{P}}
\newcommand{\E}{\mathbb{E}}
\newcommand{\eps}{\varepsilon}
\newcommand{\dd}{\,\mathrm{d}}
\begin{document}
\section*{UnsolvedMath ID: 30002813}
\textbf{Source problem code:} catalogue entry
\subsection{Definitions and mathematical statement}
In a local Euclidean interpretation, let \(n,m\) be positive integers, let \(W\subseteq\R^n\times\R^m\) be an open neighborhood of \((0,0)\), and let \(f:W\to\R^n\) be smooth with \(f(0,0)=0\). Consider the control system
\[
\dot x=f(x,u),\qquad n\notin\{2,3\}.
\]
The input dimension, admissible controls, and precise state and input domains are not specified in the supplied record; this Euclidean presentation makes the local notation explicit.

The reported hypotheses are that the Lie algebra rank condition holds at the origin and that the origin is ``locally continuously reachable in arbitrarily small time.'' The original formulation refers to these as Coron's hypotheses. The record does not define this reachability property or specify the family of vector fields used in the rank condition, so these hypotheses cannot be given an unambiguous quantified formulation from the supplied evidence alone.

For clarity, a periodic finite Fourier feedback on a state neighborhood \(U\) has the form
\[
 k(x,t)=a_0(x)+\sum_{j=1}^{N}\left[a_j(x)\cos\!\left(\frac{2\pi jt}{T}\right)+b_j(x)\sin\!\left(\frac{2\pi jt}{T}\right)\right],
\]
where \(N\in\Z_{\geq0}\), \(T>0\), and \(a_0,a_j,b_j:U\to\R^m\). Here the number of terms and the period are independent of \(x\). This is the finite Fourier interpretation of a trigonometric polynomial in time; the source does not specify the period, coefficient regularity, or whether the general formulation permits other frequency conventions.

\emph{Question as supplied.} Does every system satisfying the reported hypotheses admit a feedback \(u=k(x,t)\), trigonometric-polynomial in \(t\) with coefficients depending on \(x\), that stabilizes the origin of
\[
\dot x=f(x,k(x,t))?
\]
The original record additionally states that, in the driftless case, the target is finite Fourier feedback and asymptotic stability as specified in the report. It does not supply that specification or a definition of the driftless class. For an equilibrium target, the feedback must satisfy \(f(0,k(0,t))=0\) for all relevant \(t\); the stronger normalization \(k(0,t)=0\) is not explicitly required by the supplied text.

The general meaning of stabilization, its local or global scope, its dependence on the initial time, and the required regularity and solution concept remain unspecified. Thus this is a conditional statement of the recorded question, requiring clarification of the hypotheses and stability target before it becomes a fully precise conjecture.
\subsection{Short English statement}
For nonlinear systems satisfying the reported controllability hypotheses, with state dimension other than \(2\) or \(3\), the question asks whether stabilization can always be achieved using finitely many trigonometric terms in time whose coefficients depend on the state. The driftless version explicitly asks for asymptotic stability. Several essential definitions are missing, so the formulation requires human mathematical review.
\subsection{Sources}
\begin{itemize}
\item Local UnsolvedMath registry: supplied record with permanent ID \(30002813\), exact source problem code \texttt{catalogue entry}, title ``Trigonometric-Polynomial Feedback for Nonlinear Stabilization,'' in the Oberwolfach Reports Open Problems collection. The supplied proposed year is 2015; no individual proposer is identified.
\item Supplied report citation: \emph{Control Theory: A Mathematical Perspective on Cyber-Physical Systems} (2015), Oberwolfach Reports, DOI 10.4171/owr/2015/12, pp.\ 597--659. \url{https://doi.org/10.4171/owr/2015/12}. The record describes Oberwolfach Reports as documenting mathematical workshops and their open-problem discussions; it gives no problem-specific page or complete quotation of Coron's hypotheses.
\item Supplied literature source, labeled ``Zuyev and Grushkovskaya, systems with drift'': \url{https://arxiv.org/abs/1904.10219}. The supplied excerpt does not provide a theorem statement establishing its relationship to the full question.
\item Literature assessment checked on \texttt{2026-08-21}: the supplied curation labels the problem open and states that the extracted record does not contain a complete claim, while available constructive feedback laws cover restricted controllability classes rather than the general trigonometric-polynomial stabilization question. This is supplied dated curation derived from the OWR material, not an independently verified literature assessment; its status and cited claims require independent verification.
\item The registry labels the statement \texttt{corrected\_verified}, explaining that notation, typography, or omitted definitions were repaired while preserving mathematical content. This supplied editorial label does not establish mathematical correctness or resolve the missing specifications above. No source was independently visited for this draft.
\end{itemize}
\end{document}
